%% Đề 13 -- de-hoc-sinh-gioi-toan-9-nam-2023-2024-phong-gddt-huu-lung-lang-son.pdf
\deheader{ĐỀ SỐ 13}%
  {Kỳ thi Chọn Học sinh giỏi cấp tỉnh lớp 9 — Năm học 2023--2024}%
  {Sở GD\&ĐT Lạng Sơn | Chỉ phần Đại số}
\addcontentsline{toc}{section}{Đề 13 — Kỳ thi Chọn HSG cấp tỉnh — Năm học 2023--2024}

\begin{de_cau}{Câu 1}{4,0 điểm}
\begin{enumerate}[label=\alph*)]
  \item Rút gọn biểu thức:
  \[M = \left(\dfrac{\sqrt{a}}{\sqrt{a}-2} - \dfrac{4}{a-2\sqrt{a}}\right) : \dfrac{\sqrt{a}+2}{4}\]
  với $a > 0,\; a \neq 4$.
  \item Cho $x = \sqrt{4+2\sqrt{3}} - \sqrt{4-2\sqrt{3}}$. Chứng minh $x$ là một số nguyên.
\end{enumerate}
\end{de_cau}

\begin{loi_giai_cau}
\textbf{a)} Biến đổi:
\[\frac{\sqrt{a}}{\sqrt{a}-2} - \frac{4}{\sqrt{a}(\sqrt{a}-2)} = \frac{a-4}{\sqrt{a}(\sqrt{a}-2)} = \frac{(\sqrt{a}-2)(\sqrt{a}+2)}{\sqrt{a}(\sqrt{a}-2)} = \frac{\sqrt{a}+2}{\sqrt{a}}.\]
\[M = \frac{\sqrt{a}+2}{\sqrt{a}} \cdot \frac{4}{\sqrt{a}+2} = \boxed{\dfrac{4}{\sqrt{a}}}.\]

\textbf{b)} $\sqrt{4+2\sqrt{3}} = \sqrt{(\sqrt{3}+1)^2} = \sqrt{3}+1$.

$\sqrt{4-2\sqrt{3}} = \sqrt{(\sqrt{3}-1)^2} = \sqrt{3}-1$ (vì $\sqrt{3}>1$).

$x = (\sqrt{3}+1) - (\sqrt{3}-1) = 2 \in \mathbb{Z}$. $\square$
\end{loi_giai_cau}

\begin{de_cau}{Câu 2}{4,0 điểm}
\begin{enumerate}[label=\alph*)]
  \item Giải phương trình: $2x^2 + 3x - 5 = 0$.
  \item Giải hệ phương trình:
  \[\begin{cases} 2x + y = 5 \\ x - 3y = -1 \end{cases}.\]
\end{enumerate}
\end{de_cau}

\begin{loi_giai_cau}
\textbf{a)} $\Delta = 9 + 40 = 49$. $x = \dfrac{-3 \pm 7}{4}$: $x_1 = 1$, $x_2 = -\dfrac{5}{2}$. Vậy $\boxed{x=1}$ hoặc $\boxed{x=-\dfrac{5}{2}}$.

\textbf{b)} Từ (1): $y = 5-2x$. Thế vào (2): $x - 3(5-2x) = -1 \Rightarrow 7x = 14 \Rightarrow x=2, y=1$.

Nghiệm $\boxed{(x;y) = (2;1)}$.
\end{loi_giai_cau}

\begin{de_cau}{Câu 3}{4,0 điểm}
\begin{enumerate}[label=\alph*)]
  \item Tìm các cặp số nguyên $(x;y)$ thỏa mãn: $xy - 2x + 3y = 11$.
  \item Chứng minh rằng với mọi số lẻ $n$ thì $n^2 - 1$ chia hết cho $8$.
\end{enumerate}
\end{de_cau}

\begin{loi_giai_cau}
\textbf{a)} Biến đổi: $x(y-2)+3(y-2)=5 \Rightarrow (x+3)(y-2)=5$.

Các cặp ước: $(x+3;y-2)\in\{(1,5),(-1,-5),(5,1),(-5,-1)\}$, cho $(x;y)\in\{(-2;7),(-4;-3),(2;3),(-8;1)\}$.

\textbf{b)} Đặt $n=2k+1$ ($k\in\mathbb{Z}$): $n^2-1=4k^2+4k=4k(k+1)$.

Vì $k(k+1)$ là tích hai số nguyên liên tiếp nên chẵn, tức $k(k+1)=2m$. Suy ra $n^2-1 = 8m \vdots 8$. $\square$
\end{loi_giai_cau}

\begin{de_cau}{Câu 5}{2,0 điểm}
Cho hai số thực dương $a, b$ thỏa mãn $a + b \le 1$. Tìm giá trị nhỏ nhất của biểu thức:
\[Q = a + b + \dfrac{1}{a} + \dfrac{1}{b}.\]
\end{de_cau}

\begin{loi_giai_cau}
Áp dụng AM-GM: $a + \dfrac{1}{a} \ge 2$, $b + \dfrac{1}{b} \ge 2$, nên $Q \ge 4$. Nhưng cần xét điều kiện $a+b\le1$.

Thực ra: $Q = (a+b) + \left(\dfrac{1}{a}+\dfrac{1}{b}\right) \ge (a+b) + \dfrac{4}{a+b}$.

Đặt $s=a+b\in(0,1]$: $f(s)=s+\dfrac{4}{s} \ge 2\cdot 2 = 4$ (AM-GM, dấu $=$ khi $s=2 > 1$, không trong miền).

Trên $(0,1]$: $f'(s)=1-\dfrac{4}{s^2}<0$ nên $f$ giảm, đạt nhỏ nhất tại $s=1$: $f(1)=5$.

Vậy $Q_{\min} = \boxed{5}$ khi $a=b=\dfrac{1}{2}$.
\end{loi_giai_cau}
